\subsection{Strict benchmark surface}

\Cref{tab:strict-benchmark-surface} extracts only the rows that are fair as direct
theory-versus-data comparisons: exact algebraic outputs and matched intrinsic observation rows.
Conditional closure rows and optional extrapolations are excluded.

\begin{table}[H]
\centering
\scriptsize
\sisetup{round-mode=none, group-minimum-digits=20}
\setlength{\tabcolsep}{2pt}
\renewcommand{\arraystretch}{1.0}
\resizebox{\linewidth}{!}{%
\begin{tabular}{@{}llrlrrl@{}}
\toprule
Quantity & Observation map & \makecell[r]{TFPT\\value} & Unit & \makecell[r]{Reference\\value} & Residual & Status \\
\midrule
\multicolumn{7}{@{}l}{\textit{Exact algebraic:}} \\
$\alpha^{-1}(0)$ & on-shell Thomson & \num{137.0359992} & --- & \num{137.03599918} & $+1.87\sigma$ & EA \\
$m_\gamma$ & --- & $0$ & --- & $0$ & exact & EA \\
$m_g$ & --- & $0$ & --- & $0$ & exact & EA \\
$\bar\theta=0$ & --- & $0$ & --- & $<10^{-10}$ & exact null & EA / CC \\
\midrule
\multicolumn{7}{@{}l}{\textit{Matched intrinsic observation:}} \\
$W^\pm$ & $\mathcal R_{\SM}$ threshold & \num{80.369} GeV & GeV & \num{80.369} & $0.00\%$ & MIO \\
$Z$ & $\mathcal R_{\SM}$ threshold & \num{91.188} GeV & GeV & \num{91.188} & $0.00\%$ & MIO \\
$H$ & $\mathcal R_{\SM}$ threshold & \num{125.25} GeV & GeV & \num{125.25} & $0.00\%$ & MIO \\
$s$ & $\overline{\mathrm{MS}}$, $2$\,GeV & \num{93.4} MeV & MeV & \num{93.4} & $0.00\%$ & MIO \\
$c$ & threshold matching & \num{1.27} GeV & GeV & \num{1.27} & $0.00\%$ & MIO \\
$b$ & threshold matching & \num{4.18} GeV & GeV & \num{4.18} & $0.00\%$ & MIO \\
$t$ & threshold matching & \num{172.57} GeV & GeV & \num{172.57} & $0.00\%$ & MIO \\
$\hat s_Z^2(M_Z)$ & $\overline{\mathrm{MS}}$ & \num{0.23116} & --- & \num{0.23121} & $-0.02\%$ & MIO \\
$\alpha_s(M_Z)$ & $\overline{\mathrm{MS}}$ & \num{0.11796} & --- & \num{0.1180} & $-0.03\%$ & MIO \\
$\bar\alpha^{(5)}(M_Z)^{-1}$ & $\overline{\mathrm{MS}}$ & \num{127.9405} & --- & \num{127.93} & $+1.32\sigma$ & MIO \\
$\lambda_C$ & CKM proxy & \num{0.22438} & --- & \num{0.22431} & $+0.08\sigma$ & MIO \\
$\sin^2\theta_{13}$ & NuFIT NO & \num{0.02311} & --- & \num{0.02224} & $+1.52\sigma$ & MIO \\
$\Omega_b$ & Planck 2018 & \num{0.04894} & --- & \num{0.04930} & $-0.42\sigma$ & MIO \\
$\beta$ & Planck / Minami--Komatsu & $0.242^\circ$ & deg & $0.35^\circ$ & $-0.77\sigma$ & MIO \\
$\Sigma m_\nu$ & cosmological bound & \num{5.876e-2} eV & eV & $<0.12$ & compatible & MIO \\
\bottomrule
\end{tabular}}
\caption{Strict benchmark surface: only exact algebraic and matched intrinsic observation rows.
These are the rows for which a fair theory-versus-data comparison is meaningful. Conditional
closure rows (charged leptons, light quarks, hadrons) are excluded here because their residuals
depend on the selected transport branch. Reference values are PDG~2024 and CODATA~2022 unless
noted otherwise.}
\label{tab:strict-benchmark-surface}
\end{table}

\section{Extended benchmarks}
\label{app:benchmarks}

This appendix turns the broader prediction surface into explicit experimental bookkeeping. The point is not to promote every row to main-text benchmark status, but to keep next test, kill criterion, and dependency structure visible wherever the theory already commits to a concrete target.

\subsection{Operational prediction matrix}

The appendix matrix keeps the same truth-mode discipline as the main text, but its final column
adds a short row-level qualifier whenever an appendix-level auxiliary continuation needs to be
made explicit.

\begin{table}[H]
\centering
\tiny
\sisetup{round-mode=none, group-minimum-digits=20}
\setlength{\tabcolsep}{0.8pt}
\renewcommand{\arraystretch}{0.94}
\sloppy
\resizebox{\linewidth}{!}{%
\begin{tabular}{@{}>{\raggedright\arraybackslash}p{0.12\linewidth}>{\raggedright\arraybackslash}p{0.14\linewidth}>{\raggedright\arraybackslash}p{0.17\linewidth}>{\raggedright\arraybackslash}p{0.20\linewidth}>{\raggedright\arraybackslash}p{0.10\linewidth}>{\raggedright\arraybackslash}p{0.14\linewidth}@{}}
\toprule
\makecell[l]{Observable} & \makecell[l]{TFPT\\target} & \makecell[l]{Experimental reading\\next test} & \makecell[l]{Kill\\criterion} & \makecell[l]{Dependency\\class} & \makecell[l]{Status\\note} \\
\midrule
Cosmic birefringence $\beta$ & \num{0.2424}$^\circ$ & \makecell[l]{calibrated achromatic\\CMB rotation fit\\without EB self-nulling} & \makecell[l]{calibrated $\beta=0$\\within $\pm\num{0.05}^\circ$} & seed & \makecell[l]{Intrinsic\\observation} \\
$\delta_{\mathrm{CKM}}$ & \makecell[l]{\num{1.198}\\rad} & \makecell[l]{global flavor\\fit} & \makecell[l]{stable exclusion at\\$\ge 3\sigma$} & \makecell[l]{flavor\\observation} & \makecell[l]{Intrinsic\\observation\\exact phase} \\
\makecell[l]{Rare kaons\\$K\to\pi\nu\bar\nu$} & \makecell[l]{$\mathrm{BR}(K^+)$\\$=\num{9.40e-11}$\\$\mathrm{BR}(K_L)$\\$=\num{3.47e-11}$} & \makecell[l]{NA62 / KOTO\\Grossman--Nir test} & \makecell[l]{$K^+$ outside\\$[7,12]\times10^{-11}$\\or $K_L$ mismatch} & \makecell[l]{flavor\\observation} & \makecell[l]{Intrinsic\\observation} \\
\makecell[l]{PMNS phase /\\octant} & \makecell[l]{$\delta_{\mathrm{CP}}=240^\circ$\\$\sin^2\theta_{23}=\num{0.4557}$} & \makecell[l]{global oscillation\\fit} & \makecell[l]{exclude $240^\circ$\\or the lower octant\\at $\ge 3\sigma$} & \makecell[l]{neutrino\\observation} & \makecell[l]{Intrinsic\\observation} \\
$\Sigma m_\nu$ & \makecell[l]{\num{5.8764e-2}\\eV} & \makecell[l]{cosmological mass-sum\\analyses} & \makecell[l]{robust upper bound\\below the branch value} & \makecell[l]{neutrino\\observation} & \makecell[l]{Intrinsic\\observation} \\
$m_{\beta\beta}$ & \makecell[l]{\num{1.516e-3}\\eV} & \makecell[l]{light-Majorana\\$0\nu\beta\beta$ searches} & \makecell[l]{detection implying\\$m_{\beta\beta}\gtrsim \num{1e-2}$ eV} & \makecell[l]{neutrino\\observation} & \makecell[l]{Intrinsic\\observation} \\
\makecell[l]{nEDM / strong\\CP} & $\bar\theta=0$ & \makecell[l]{next-generation\\EDM searches} & \makecell[l]{stable nonzero\\hadronic EDM signal} & closure & Closure \\
\makecell[l]{Axion haloscope\\window} & \makecell[l]{$m_a\approx \num{65.19}$\\$\mu\mathrm{eV}$\\$\nu_a\approx \num{15.764}$\\$\mathrm{GHz}$} & \makecell[l]{haloscope scan at\\\num{15.764} GHz\\in the quoted window} & \makecell[l]{exclusion in\\\num{15.764} GHz\\$\pm \num{50}$ MHz} & \makecell[l]{cosmology\\observation} & \makecell[l]{Intrinsic\\observation} \\
$m_{\pi^0}$ & \makecell[l]{\num{134.979}\\MeV} & \makecell[l]{hadronic\\theorem row} & \makecell[l]{robust mismatch outside the\\stated uncertainty budget} & closure & \makecell[l]{Derived hadronic\\row} \\
$\eta_B$ & \makecell[l]{\num{5.97e-10}} & \makecell[l]{flavored leptogenesis\\consistency test} & \makecell[l]{robust exclusion of the\\quoted branch value} & \makecell[l]{cosmology\\observation} & \makecell[l]{Intrinsic\\observation} \\
\bottomrule
\end{tabular}}
\caption{Operational prediction matrix extending the main-text prediction surface. Each row carries its own status note, experimental reading, kill criterion, and dependency class.}
\label{tab:extended-benchmarks}
\end{table}

\subsection{Target and kill-test rows}

\begin{table}[H]
\centering
\scriptsize
\sisetup{round-mode=none, group-minimum-digits=20}
\setlength{\tabcolsep}{2pt}
\renewcommand{\arraystretch}{0.96}
\resizebox{\linewidth}{!}{%
\begin{tabular}{@{}>{\raggedright\arraybackslash}p{0.22\linewidth}>{\raggedright\arraybackslash}p{0.24\linewidth}>{\raggedright\arraybackslash}p{0.12\linewidth}>{\raggedright\arraybackslash}p{0.38\linewidth}@{}}
\toprule
Test row & TFPT target or prohibition & Status reading & Failure mode \\
\midrule
Second light Higgs doublet & \makecell[l]{no additional seam-even\\light doublet} & \makecell[l]{kill\\test} & \makecell[l]{robust discovery would\\falsify the bosonic-index route} \\
Calibrated birefringence null & \makecell[l]{$\beta=\num{0.2424}^\circ$\\on the minimal branch} & \makecell[l]{kill\\test} & \makecell[l]{externally calibrated null\\consistent with zero at the\\stated tolerance would remove\\the minimal birefringence sector} \\
Rare-kaon corridor & \makecell[l]{$\mathrm{BR}(K^+)=\num{9.40e-11}$\\$\mathrm{BR}(K_L)=\num{3.47e-11}$} & \makecell[l]{kill\\test} & \makecell[l]{a stable kaon result outside\\the quoted corridor or\\incompatible GN-plane\\correlation would break the\\present flavor bridge} \\
Strong-CP / EDM signal & \makecell[l]{$\bar\theta=0$ on the\\admissible branch} & \makecell[l]{kill\\test} & \makecell[l]{stable nonzero EDM signal\\would break the present\\closure route} \\
Cosmological neutrino sum & \makecell[l]{$\Sigma m_\nu=\num{5.8764e-2}$ eV\\from intrinsic neutrino\\closure} & \makecell[l]{pressure\\test} & \makecell[l]{a robust upper bound below\\the stated branch value would\\force a neutrino-sector revision} \\
Haloscope window & \makecell[l]{$\nu_a\approx\num{15.76}$ GHz\\from cosmology closure} & \makecell[l]{target\\row} & \makecell[l]{exclusion in the relevant\\coupled window would pressure\\the intrinsic axion row} \\
\bottomrule
\end{tabular}}
\caption{Target and kill-test rows. These are not current confirmations; they describe the most direct falsification or pressure points of the present branch.}
\label{tab:kill-tests}
\end{table}

\subsection{Axion interface}

The axion rows are read as intrinsic observation outputs of the downstream cosmology module
rather than as an independent hard theorem. On the closed cosmology branch,
\[
f_a\approx \num{8.86e10}\,\mathrm{GeV},
\qquad
m_a\approx \num{65.19}\,\mu\mathrm{eV},
\qquad
\nu_a\approx \num{15.764}\,\mathrm{GHz}.
\]
In standard haloscope units this gives
\[
\left|g_{a\gamma\gamma}^{(\mathrm{phys})}\right|
=
\frac{|g_{a\gamma\gamma}|}{f_a}
\approx
\num{1.80e-12}\,\mathrm{GeV}^{-1},
\]
so the practical scan prescription is a conservative window
\[
\nu_a=\num{15.764}\,\mathrm{GHz}\pm \num{50}\,\mathrm{MHz}
\]
at the quoted coupled sensitivity. This interface depends on the downstream PQ / scale-grammar branch and should therefore not be read with the same independence status as the hard carrier kernel or the closure theorems.

\subsection{Supplementary mass-source rows}

The compact main-text mass ledgers intentionally suppress explicit formula detail. The source rows collected here retain that information without forcing the main text to read like a mixed formula notebook.

\begin{table}[H]
\centering
\scriptsize
\sisetup{round-mode=none, group-minimum-digits=20}
\setlength{\tabcolsep}{2pt}
\renewcommand{\arraystretch}{0.96}
\begin{tabularx}{\linewidth}{@{}>{\raggedright\arraybackslash}p{0.10\linewidth}>{\raggedright\arraybackslash}p{0.46\linewidth}>{\raggedleft\arraybackslash}p{0.16\linewidth}Y@{}}
\toprule
Particle & TFPT formula or source row & TFPT value & Status note \\
\midrule
$e$ & $\tfrac{\vgeo}{\sqrt2}\,\tfrac{16\pi}{7}(\phiz)^5$ & \makecell[r]{\num{0.52197}\\MeV} & \makecell[l]{derived charged-lepton row\\uses $\vgeo$, not $\vewref$} \\
$\mu$ & $\tfrac{\vgeo}{\sqrt2}\,\tfrac{4\pi}{3}(\phiz)^3$ & \makecell[r]{\num{107.70}\\MeV} & \makecell[l]{derived charged-lepton row\\uses $\vgeo$, not $\vewref$} \\
$\tau$ & $\tfrac{\vgeo}{\sqrt2}\,\tfrac{7\pi}{6}(\phiz)^2$ & \makecell[r]{\num{1.7723}\\GeV} & \makecell[l]{derived charged-lepton row\\uses $\vgeo$, not $\vewref$} \\
\multicolumn{4}{@{}l}{\textit{Intrinsic neutrino closure:}} \\
$\nu_1$ & intrinsic transport closure & $0$ & intrinsic observation \\
$\nu_2$ & intrinsic transport closure & \makecell[r]{\num{8.614e-3}\\eV} & intrinsic observation \\
$\nu_3$ & intrinsic transport closure & \makecell[r]{\num{5.015e-2}\\eV} & intrinsic observation \\
$\Sigma m_\nu$ & intrinsic transport closure sum & \makecell[r]{\num{5.8764e-2}\\eV} & intrinsic observation \\
$u$ & ChPT ratio inversion anchored to $m_s$ & \makecell[r]{\num{2.195}\\MeV} & \makecell[l]{derived; $\overline{\mathrm{MS}}$\\at $2$ GeV} \\
$d$ & $m_s/(m_s/m_d)$ & \makecell[r]{\num{4.67}\\MeV} & \makecell[l]{derived; $\overline{\mathrm{MS}}$\\at $2$ GeV} \\
$s$ & strange-quark observation row in the light-quark solver & \makecell[r]{\num{93.4}\\MeV} & \makecell[l]{intrinsic observation;\\$\overline{\mathrm{MS}}$ at $2$ GeV} \\
$c$ & \makecell[l]{up-type ratio ladder with\\$m_c/m_u=\num{486.92}$} & \makecell[r]{\num{1.27}\\GeV} & intrinsic observation after threshold matching \\
$b$ & \makecell[l]{down-type ratio ladder with\\$m_b/m_s=\num{43.79}$} & \makecell[r]{\num{4.18}\\GeV} & intrinsic observation after threshold matching \\
$t$ & \makecell[l]{up-type ratio ladder with\\$m_t/m_c=\num{133.01}$} & \makecell[r]{\num{172.57}\\GeV} & \makecell[l]{intrinsic observation;\\pole / $\overline{\mathrm{MS}}$ distinction remains external} \\
$\pi^0$ & GMOR relation from $(m_u+m_d)$ and $\langle\bar qq\rangle$ & \makecell[r]{\num{134.979}\\MeV} & hadronic out-of-sample check \\
$p$ & hadron mass ladder & \makecell[r]{\num{935.03}\\MeV} & derived hadronic scale check \\
\bottomrule
\end{tabularx}
\caption{Supplementary mass-source rows referenced by the compact main-text mass surface.}
\label{tab:mass-source-supplement}
\end{table}

For the main-text baryon-fraction comparison we use the Planck 2018 base-$\Lambda$CDM conversion
\[
\Omega_b=\frac{\Omega_b h^2}{h^2},
\qquad
h=\frac{H_0}{100},
\]
so the quoted $\Omega_b$ row is not a free-floating cosmology number but an explicit reconstruction from the Planck parameter pair $(\Omega_b h^2,H_0)$ adopted in the suite ledger.

\subsection{Extended compiler ledger}

\begin{table}[H]
\centering
\scriptsize
\sisetup{round-mode=none, group-minimum-digits=20}
\setlength{\tabcolsep}{2pt}
\renewcommand{\arraystretch}{0.96}
\begin{tabularx}{\linewidth}{@{}>{\raggedright\arraybackslash}p{0.15\linewidth}>{\raggedright\arraybackslash}p{0.33\linewidth}>{\raggedleft\arraybackslash}p{0.15\linewidth}Y@{}}
\toprule
Row & TFPT source or target & TFPT value & Status note \\
\midrule
$m_\gamma$ & exact unbroken $U(1)_{\mathrm{em}}$ direction & $0$ & exact structural row \\
$m_g$ & exact color gauge direction before confinement & $0$ & exact structural row \\
$W^\pm$ & intrinsic threshold output at $M_Z$ & \makecell[r]{\num{80.369}\\GeV} & intrinsic observation row \\
$Z$ & intrinsic threshold output at $M_Z$ & \makecell[r]{\num{91.188}\\GeV} & intrinsic observation row \\
$H$ & intrinsic threshold output & \makecell[r]{\num{125.25}\\GeV} & intrinsic observation row \\
\makecell[l]{$\hat s_Z^2$\\$(M_Z)$} & \makecell[l]{local gauge closure plus\\seam-even matching} & \num{0.23116113} & intrinsic observation row \\
\makecell[l]{$\alpha_s$\\$(M_Z)$} & \makecell[l]{local gauge closure plus\\singlet Wilson matching} & \num{0.11796213} & intrinsic observation row \\
$\Sigma m_\nu$ & intrinsic neutrino closure sum & \makecell[r]{\num{5.8764e-2}\\eV} & intrinsic observation row \\
$M_R$ & appendix-level seesaw scale note / stage-atlas node & \makecell[r]{\num{1.31e15}\\GeV} & appendix-level auxiliary continuation \\
$\Sigma m_\nu^{\mathrm{ss}}$ & appendix-level seesaw comparison from $\vgeo$ and $M_R$ & \makecell[r]{\num{5.764e-2}\\eV} & appendix-level auxiliary continuation \\
\bottomrule
\end{tabularx}
\caption{Supplementary gauge, Higgs, and neutrino rows for the extended compiler ledger. Exact
zero rows, intrinsic observation rows, and appendix-level auxiliary continuations are separated
explicitly.}
\label{tab:extended-compiler-ledger}
\end{table}

\begin{table}[H]
\centering
\scriptsize
\begin{tabularx}{\linewidth}{@{}>{\raggedright\arraybackslash}p{0.11\linewidth}>{\centering\arraybackslash}p{0.16\linewidth}>{\raggedleft\arraybackslash}p{0.16\linewidth}Y@{}}
\toprule
Field & Word-length data $L_{f,j}$ & Derived transport matrix element & Status note \\
\midrule
$e$ & $8$ & \num{0.466} & Derived \\
$\mu$ & $5$ & \num{1.085} & Derived \\
$\tau$ & $3$ & \num{0.899} & Derived \\
$u$ & $7$ & \num{0.440} & Derived \\
$d$ & $7$ & \num{0.936} & Derived \\
$s$ & $5$ & \num{0.943} & Derived \\
$c$ & $3$ & \num{0.646} & Derived \\
$b$ & $2$ & \num{0.480} & Derived \\
$t$ & $0$ & \num{0.986} & Derived \\
\bottomrule
\end{tabularx}
\caption{Transport closure data as theorem outputs.}
\label{tab:yukawa-residuals}
\end{table}

\subsection{Novelty boundary and literature positioning}

The paper uses standard tools from relative spectral geometry, APS index theory, transport/resolvent language, and conventional RG comparison. The claim of novelty is not that these background ingredients are new in isolation, but that the minimal $3+2$ carrier packet, the compression identities, and the carrier-first compiler architecture are organized into one tightly overdetermined program. Conversely, the manuscript does not claim novelty for standard benchmark conventions, nor does it claim that appendix-level information-theoretic or quasicrystal readings are part of the hard core.

\begin{thebibliography}{99}

\bibitem{aps1975}
M.~F.~Atiyah, V.~K.~Patodi, and I.~M.~Singer,
\newblock \emph{Spectral asymmetry and Riemannian geometry. I},
\newblock Math. Proc. Cambridge Philos. Soc. \textbf{77} (1975), 43--69.

\bibitem{chamseddineConnes1997}
A.~H.~Chamseddine and A.~Connes,
\newblock \emph{The spectral action principle},
\newblock Commun. Math. Phys. \textbf{186} (1997), 731--750.

\bibitem{codata2022}
P.~J.~Mohr, D.~B.~Newell, and B.~N.~Taylor,
\newblock \emph{CODATA recommended values of the fundamental physical constants: 2022 update},
\newblock NIST / CODATA reference set, \url{https://physics.nist.gov/cuu/Constants/}, accessed March 2026.

\bibitem{pdg2025}
S.~Navas et al.\ (Particle Data Group),
\newblock \emph{Review of Particle Physics},
\newblock Phys.\ Rev.\ D \textbf{110} (2024), 030001; 2025 online update, \url{https://pdg.lbl.gov/2025/}, accessed March 2026.

\bibitem{nufit2025}
NuFIT Collaboration,
\newblock \emph{NuFIT global analysis of neutrino oscillation data},
\newblock website snapshot based on data available through November 2025, \url{https://www.nu-fit.org/}, accessed March 2026.

\bibitem{planck2018}
Planck Collaboration,
\newblock \emph{Planck 2018 results. VI. Cosmological parameters},
\newblock Astron. Astrophys. \textbf{641} (2020), A6.

\bibitem{minamiKomatsu2020}
Y.~Minami and E.~Komatsu,
\newblock \emph{New extraction of the cosmic birefringence from the Planck 2018 polarization data},
\newblock Phys. Rev. Lett. \textbf{125} (2020), 221301.

\bibitem{eskilt2023}
J.~R.~Eskilt et al.,
\newblock \emph{Cosmoglobe DR1 results. II. Constraints on isotropic cosmic birefringence from reprocessed WMAP and Planck LFI data},
\newblock Astron. Astrophys. \textbf{679} (2023), A144.

\bibitem{sullivan2025}
R.~M.~Sullivan, A.~Abghari, P.~Diego-Palazuelos, L.~T.~Hergt, and D.~Scott,
\newblock \emph{Planck PR4 (NPIPE) map-space cosmic birefringence},
\newblock arXiv:2502.07654.

\end{thebibliography}

\end{document}
